
Explainable Artificial Intelligence (XAI) refers to 
a set of methods and techniques that make the 
predictions of machine learning models interpretable 
and transparent. In the context of this thesis, XAI 
methods are employed to reveal and interpret memorized 
features in CNNs trained on image classification tasks. 
Formally, given a model $f_\theta$ and an input image 
$x$, an attribution method computes a saliency map 
$\mathcal{A}(x) \in \mathbb{R}^{H \times W}$ that 
assigns an importance score to each pixel, indicating 
its contribution to the model's prediction.

\subsubsection{Gradient-weighted Class Activation 
Mapping (GradCAM)}

GradCAM (Selvaraju et al., 2017) computes the 
importance of each feature map in the final 
convolutional layer with respect to a target 
class $c$. The weights $\alpha_k^c$ for feature 
map $k$ are obtained by global average pooling 
of the gradients:

\begin{equation}
\alpha_k^c = \frac{1}{Z} \sum_{i} \sum_{j} 
\frac{\partial \hat{y}^c}{\partial A^k_{ij}}
\label{eq:gradcamweights}
\end{equation}

where $Z$ is the number of pixels in the feature 
map, $\hat{y}^c$ is the predicted score for class 
$c$, and $A^k_{ij}$ is the activation of feature 
map $k$ at spatial location $(i,j)$. The GradCAM 
saliency map is then computed as:

\begin{equation}
\mathcal{L}^c_{GradCAM} = \text{ReLU}
\left(\sum_{k} \alpha_k^c A^k\right)
\label{eq:gradcam}
\end{equation}

The ReLU operation retains only the features that 
have a positive influence on the target class $c$. 
GradCAM produces a coarse localization map that 
highlights the regions of the input image most 
relevant to the model's prediction, making it 
particularly suitable for revealing memorized 
spatial features in CNNs.

\subsubsection{Integrated Gradients}

Integrated Gradients (Sundararajan et al., 2017) 
computes attributions by integrating the gradients 
of the model's output with respect to the input 
along a straight path from a baseline $x'$ 
(typically a black image) to the actual input $x$:

\begin{equation}
\mathcal{IG}_i(x) = (x_i - x'_i) \times 
\int_{\alpha=0}^{1} \frac{\partial f_\theta
(x' + \alpha(x - x'))}{\partial x_i} d\alpha
\label{eq:ig}
\end{equation}

where $x_i$ is the $i$-th input feature, $x'_i$ 
is the corresponding baseline feature, and 
$\alpha$ interpolates between the baseline and 
the input. In practice, the integral is 
approximated using a Riemann sum with $m$ steps:

\begin{equation}
\mathcal{IG}_i(x) \approx (x_i - x'_i) \times 
\sum_{k=1}^{m} \frac{\partial f_\theta\left(x' + 
\frac{k}{m}(x - x')\right)}{\partial x_i} 
\times \frac{1}{m}
\label{eq:igapprox}
\end{equation}

Integrated Gradients satisfies the completeness 
axiom, meaning that the sum of attributions 
equals the difference between the model's output 
at the input and the baseline:

\begin{equation}
\sum_{i} \mathcal{IG}_i(x) = 
f_\theta(x) - f_\theta(x')
\label{eq:igcompleteness}
\end{equation}

\subsubsection{SHAP (SHapley Additive 
exPlanations)}

SHAP (Lundberg \& Lee, 2017) is grounded in 
cooperative game theory and computes the 
contribution of each feature to the model's 
prediction based on Shapley values. For a 
model $f_\theta$ and input $x$, the SHAP 
value for feature $i$ is defined as:

\begin{equation}
\phi_i = \sum_{S \subseteq F \setminus \{i\}} 
\frac{|S|!(|F|-|S|-1)!}{|F|!} 
\left[f_\theta(S \cup \{i\}) - f_\theta(S)\right]
\label{eq:shap}
\end{equation}

where $F$ is the set of all features, $S$ is a 
subset of features excluding $i$, and 
$f_\theta(S)$ denotes the model's prediction 
using only the features in $S$. The Shapley 
value $\phi_i$ represents the average marginal 
contribution of feature $i$ across all possible 
feature subsets, ensuring a fair and consistent 
attribution.

\subsubsection{LIME (Local Interpretable 
Model-agnostic Explanations)}

LIME (Ribeiro et al., 2016) approximates the 
behaviour of a complex model locally around 
a specific input $x$ by fitting an interpretable 
surrogate model $g \in G$:

\begin{equation}
\xi(x) = \arg\min_{g \in G} 
\mathcal{L}(f_\theta, g, \pi_x) + \Omega(g)
\label{eq:lime}
\end{equation}

where $\mathcal{L}(f_\theta, g, \pi_x)$ measures 
how well the surrogate model $g$ approximates 
$f_\theta$ in the neighborhood of $x$ defined 
by the proximity measure $\pi_x$, and $\Omega(g)$ 
is a complexity penalty that ensures the surrogate 
model remains interpretable. For image 
classification, LIME generates perturbed versions 
of the input by masking superpixels and observes 
the change in the model's prediction to identify 
the most influential regions.

\subsubsection{Connecting XAI to Memorization}

The attribution methods described above serve 
as diagnostic tools for revealing memorization 
in CNNs. A model that has memorized specific 
training instances will produce saliency maps 
that highlight instance-specific spurious 
features rather than semantically meaningful 
regions. Formally, for a memorized example 
$(x_i, y_i)$ with high memorization score 
$\text{mem}(f, \mathcal{D}, i)$ as defined 
in Equation \ref{eq:memscore}, the attribution 
map $\mathcal{A}(x_i)$ is expected to assign 
high importance to spurious features $v_s$ 
rather than core features $v_c$:

\begin{equation}
\sum_{j \in v_s} \mathcal{A}_j(x_i) \gg 
\sum_{j \in v_c} \mathcal{A}_j(x_i)
\label{eq:xaimem}
\end{equation}

This formal connection between memorization 
and XAI attributions constitutes the 
theoretical foundation of the audit pipeline 
proposed in Section 3.5.